% ===== tr-06 块2：第 6 章 pp.13--16（原文页 13--16）=====
% [边界] 本文件从 p.13 页首（第 5 节标题）起。p.12 末行为脚注 4，与 p.13 正文不构成跨页句；p.13 由第 5 节标题开始，故本文件直接从 \section 起笔。
% [原文 p.13]
\section{超 Kähler 度量的唯一性定理}

先验地，Hein 的构造 [32] 或其在定理 6.4.2 中的推广，可能在同一 de Rham 上同调类、甚至同一 Bott--Chern 上同调类中产生许多超 Kähler 度量。另一方面，从 Calabi 猜想的原始证明 [64] 出发，人们预期复 Monge--Ampère 方程在适当指定的渐近行为下具有唯一解。

即使有 Hein 的强估计给出超 Kähler 度量到模型度量的快速衰减，建立唯一性定理也并非易事。主要障碍在于 $\check X$ 不是 Stein 的，且容许许多全纯向量场。事实上，由于 $\check X$ 不是 Stein 的，同一 de Rham 上同调类中的两个给定超 Kähler 度量可能不代表同一个 Bott--Chern 上同调类。本节我们概述 [11] 中建立的超 Kähler 度量的唯一性定理。粗略地说，$\check X$ 上 $\sqrt{-1}\partial\bar\partial$-引理的失效可以直接与 $\check X$ 上的全纯向量场联系起来。

\begin{theorem}{6.5.1}[{[11, 定理 1.4]}]
设 $\check\omega_1,\check\omega_2$ 是 $\check X$ 上的两个超 Kähler 度量，满足：
\begin{enumerate}
\item $[\check\omega_1]_{dR}=[\check\omega_2]_{dR}\in H^2(\check X,\R)$；
\item $\check\omega_i^2=\alpha^2\Omega\wedge\bar\Omega$，其中 $\alpha>0$，$i=1,2$；
\item 存在一个（可能非标准的）半平坦度量 $\check\omega_{sf,b_0,\varepsilon}$，使得在无穷远处适当选取局部截面之后，有
\begin{equation}
|\check\omega_i-\alpha\check\omega_{sf,\sigma,b_0,\varepsilon}|_{\check\omega_{sf,b_0,\varepsilon}}\le Cr^{-\frac43}. \tag{6.5.1}
\end{equation}
\end{enumerate}
则 $\check\omega_1=\check\omega_2$。
\end{theorem}
% [Tutor Note: 本节定理断言中的度量都带 check 帽（$\check\omega_i$、$\check X$），但证明部分原书把帽子几乎全部省略、写成 $\omega_i$（例如“given $\omega_1,\omega_2$ as above”、“$B_{\omega_i}(R)$”）；页图上可逐处核对，勿误以为证明了另一个对象。另外 (5.1) 中第一个度量的下标多出 $\sigma$（截面），而作为范数的度量下标为 $sf,b_0,\varepsilon$，页图如此，照录。]

\begin{proof}[证明概要]
设给定如上 $\omega_1,\omega_2$。取定点 $x_0\in X$，以 $B_{\omega_i}(R)$ 记 $x_0$ 周围半径为 $R$ 的球。由于 $\omega_i$ 收敛到半平坦度量，我们有
\[
\operatorname{Vol}(B_{\omega_i}(R))\sim R^{4/3}.
\]
设我们可把 $\omega_2$ 写成 $\omega_2=\omega_1+\sqrt{-1}\partial\bar\partial\varphi$，其中 $\varphi$ 有界。通过加常数，可设 $\varphi\ge0$。于是由 Monge--Ampère 方程与 AM--GM 不等式，我们有
\[
\Delta_{\omega_1}\varphi\ge0.
\]
令 $r(x)=\operatorname{dist}_{\omega_1}(x_0,x)$，并取光滑函数 $\eta\ge0$，使得 $\eta\equiv1$ 在 $[0,1]$ 上、$\eta\equiv0$ 在 $[2,\infty]$ 上，且 $|\eta'|\le2$。定义
\[
\eta_R=\eta\left(\frac rR\right).
\]
则对 $R\gg1$ 有
\begin{align*}
\int_X\eta_R^2\Delta_{\omega_1}\varphi^2\omega_1^2
&=2\int_X\eta_R^2\bigl(|\nabla\varphi|^2_{\omega_1}+\varphi\Delta_{\omega_1}\varphi\bigr)\omega_1^2\\
&\ge2\int_X\eta_R^2|\nabla\varphi|^2_{\omega_1}\omega_1^2.
\end{align*}
% [Tutor Note: 左端的 $\Delta_{\omega_1}\varphi^2$ 应读作 Laplacian 作用在函数 $\varphi^2$ 上，而不是 $(\Delta_{\omega_1}\varphi)^2$：右端恰是点态恒等式 $\Delta(\varphi^2)=2(|\nabla\varphi|^2+\varphi\Delta\varphi)$ 乘以 $\eta_R^2\omega_1^2$ 后逐点相等；若读成 $(\Delta\varphi)^2$ 则等号不成立。]
% [原文 p.14]
另一方面，分部积分给出
\begin{align*}
\int_X\eta_R^2\Delta_{\omega_1}\varphi^2\omega_1^2
&=-4\int_X\eta_R\varphi\langle\nabla\eta_R,\nabla\varphi\rangle_{\omega_1}\omega_1^2\\
&\le\int_X\eta_R^2|\nabla\varphi|^2_{\omega_1}\omega_1^2+C\int_X\varphi^2|\nabla\eta_R|^2_{\omega_1}\omega_1^2.
\end{align*}

结合这些不等式，并利用 $\varphi$ 有界以及 $|\nabla\eta_R|\le2R^{-1}$，可得
\[
\int_X\eta_R^2|\nabla\varphi|^2_{\omega_1}\omega_1^2\le CR^{-2}\int_{B(2R)}\omega_1^2\le CR^{4/3-2}.
\]
令 $R\to+\infty$ 取极限，得 $\nabla\varphi=0$，从而 $\varphi$ 为常数。于是，我们把唯一性定理归结为构造一个有界位势函数 $\varphi$。

为构造有界位势，我们如下论证。写 $\omega_2=\omega_1+d\beta$。设 $\omega_{sf,b_0,\varepsilon}$ 在无穷远处由某个截面 $\sigma$ 诱导。我们可以找到一个紧化 $X$ 的有理椭圆曲面 $Y$，使得 $\sigma$ 延拓为 $I_d$ 奇异纤维 $Y\setminus X$ 上的全纯截面。按假设我们有
\begin{equation}
|d\beta|_{\omega_{sf,b_0,\frac{\varepsilon}{\alpha}}}\le Cr^{-4/3}. \tag{6.5.2}
\end{equation}
% [Tutor Note: (5.2) 的范数下标是 $\omega_{sf,b_0,\varepsilon/\alpha}$（参数已按 $\alpha$ 归一化，页图上 $\varepsilon$ 与 $\alpha$ 排成上下相叠的小分式），不是 $\omega_{sf,b_0,\varepsilon}$ 再乘一个因子 $\alpha$；与 (5.3) 的写法一致。]

通过把半平坦度量与 $I_d$ 纤维附近坐标中的光滑 Kähler 度量作比较，可以看到 $d\beta$ 具有 $L^1$ 系数，从而延拓为 $Y$ 上的一个 $(1,1)$-流 $T$。接着可以证明 $T$ 是闭的，且 $[T]_{dR}=0\in H^2(Y,\R)$。于是由流的 $\sqrt{-1}\partial\bar\partial$-引理，可写 $T=\sqrt{-1}\partial\bar\partial\varphi$，其中 $\varphi$ 是 $L^1$ 函数。现在需要证明 $\varphi$ 有界。不妨像前面那样在 $\check X_{\mathrm{mod}}$ 中用坐标 $(u,v)$ 工作。结合 $d\beta=\sqrt{-1}\partial\bar\partial\varphi$ 的界 (6.5.2)，可推出沿 $\check X_{\mathrm{mod}}\to\Delta^*$ 的平坦环面纤维有 $|\varphi_{v\bar v}|\le C/(-\log|u|)^3$，所以 $\varphi$ 沿纤维大致是调和的。直观上，这意味着 $\varphi$ 沿纤维大致为常数。确实，由于纤维是平坦环面，用椭圆正则性以及 Sobolev 与 Poincaré 不等式作论证可得
\[
\sup_{\pi^{-1}(u)}\left|\varphi-\frac{2\pi}{d(-\log|u|)}\int_{\pi^{-1}(u)}\varphi\,\frac{\sqrt{-1}\,dv\wedge d\bar v}{2}\right|\le\frac{C}{(-\log|u|)^{\frac12}}.
\]
粗略地说，这意味着我们可以假定 $\varphi=\varphi(u)$ 是从椭圆纤维化的底空间拉回的某个函数。在这种情形下，估计 (6.5.2) 给出
\[
|\varphi_{u\bar u}|\le\frac{C}{|u|^2(-\log|u|)^3},
\]
它容易推出 $u$ 可以延拓为 $I_d$ 纤维上的有界函数。\proofend
\end{proof}
% [Tutor Note: 末句原作“which easily implies that u can be extended over the $I_d$ fiber as a bounded function”（页 14，x 60--570pt、y 630--650pt 区域 800dpi 可辨），按上下文能延拓的应是被估计的函数 $\varphi=\varphi(u)$ 而非坐标 $u$ 本身，疑为原文笔误；译文照录原样，此注给出正确理解。]

对后续应用而言，关键点是：我们只假设超 Kähler 度量 $\check\omega_i$ 到半平坦度量有多项式衰减，而不是指数衰减。事实上，给定一个可能非标准的半平坦
% [原文 p.15]
度量，人们可以研究沿某个截面作平移对其渐近性与 Bott--Chern 上同调类的影响。利用这一点，可以从定理 6.5.1 推出如下更强的结果：

\begin{theorem}{6.5.2}[{[11, 命题 4.8]}]
以 $\operatorname{Aut}_0(\check X)$ 记 $\check X$ 的、同伦于恒等的保纤维双全纯映射。设 $\check\omega_1,\check\omega_2$ 是 $\check X$ 上两个完全 Ricci 平坦度量，满足
\begin{enumerate}
\item $\check\omega_1^2=\check\omega_2^2$；
\item $[\check\omega_1]_{dR}=[\check\omega_2]_{dR}$；
\item 存在（可能非标准的）半平坦度量 $\omega_{sf,\sigma_i,b_{0,i},\varepsilon_i}$，使得
\[
|\check\omega_i-\alpha\check\omega_{sf,b_{0,i},\frac{\varepsilon_i}{\alpha}}|_{\check\omega_i}\le Cr_i^{-\frac43},
\]
其中 $r_i$ 是相对于 $\check\omega_i$ 到一固定点的距离。
\end{enumerate}
则存在 $\Phi\in\operatorname{Aut}_0(\check X)$ 使得 $\Phi^*\check\omega_1=\check\omega_2$。
\end{theorem}

为证明定理 6.5.2，思路是造一个整体截面 $\tau:\C\to\check X$，使得与沿纤维平移 $\tau$ 相对应的映射 $\Phi$ 与恒等映射同痕，并且
\begin{equation}
|\Phi^*\check\omega_{sf,b_{0,2},\frac{\varepsilon_i}{\alpha}}-\check\omega_{sf,b_{0,1},\frac{\varepsilon_i}{\alpha}}|_{\check\omega_{sf,b_{0,1},\frac{\varepsilon_i}{\alpha}}}\le Cr^{-\frac43}. \tag{6.5.3}
\end{equation}
于是便可应用定理 6.5.1 得出结论。只要求多项式衰减而非指数衰减，其重要性在于：可能不存在能给出 式 (6.5.3) 中指数衰减的整体自同构 $\Phi$。

上述定理启发了下列关于 $\check X$ 的 Kähler 模空间的定义，它为第 6.6.1 节的 SYZ 镜像对称讨论奠定了基础。固定全纯体积形式 $\check\Omega$，把 $\check X$ 的完全、渐近半平坦的超 Kähler 度量的模空间定义为
\[
\check K_{CY}:=\{\text{Kähler 度量 }\check\omega\mid 2\check\omega^2=\check\Omega\wedge\bar{\check\Omega}\text{ 且满足式 (6.5.1)}\}/\operatorname{Aut}_0(\check X).
\]
由定理 6.5.2，我们可以把 $\check K_{CY}$ 等同于 $H^2(\check X,\R)$ 的一个子集。继而可以证明 $\check K_{CY}$ 是 $H^2(\check X,\R)$ 的开子集。特别地，若 $\check D$ 是一个 $I_d$-纤维，则 $\check K_{CY}$ 的维数为 $11-d$。注意，由次数为 $d$ 的 del Pezzo 曲面与光滑反典范除子构成的配对 $(Y,D)$ 的模空间维数为 $10-d$。特别地，$\check K_{CY}$ 似乎多出了一维。

通过限制到 $H^2(\check X,\R)$ 中容许特殊拉格朗日纤维化的 Kähler 类的子集，这一不匹配得到了修正。

\section{应用}

\subsection{对数 Calabi--Yau 曲面的 SYZ 镜像对称}

众所周知，del Pezzo 曲面 $Y$ 的镜像是 Landau--Ginzburg 模型，即一个全纯函数 $W:(\C^*)^2\to\C$，称为超势。超势 $W$ 在多种意义下是 $Y$ 的镜像：
% [原文 p.16]
\begin{itemize}
\item 在计数几何层面，$W$ 的量子周期计算 $Y$ 的后代 Gromov--Witten 不变量 [22][18][8]。
\item 在同调镜像对称层面，猜想 $X$ 的 Fukaya 范畴等价于 $W$ 的奇点导出范畴或 $W$ 的矩阵分解范畴 [21][48][9]。反过来，$W$ 的 Fukaya--Seidel 范畴等价于 $Y$ 上凝聚层的导出范畴 [2]。特别地，Auroux--Katzarkov--Orlov [2] 观察到：若 $Y$ 的次数为 $d$，则通过添加一个 $I_d$ 纤维，$W$ 可以在拓扑上紧化为一个有理椭圆曲面。
\item Lunts--Przyjalkowski [51] 证明了双方在 Katzarkov--Kontsevich--Pantev [43] 意义下的 Hodge 数相同。
\item Doran--Thompson [17] 提出了带光滑反典范除子的次数 $d$ 的 del Pezzo 曲面配对与带 $I_d$-纤维的有理椭圆曲面之间的格极化镜像对称。
\end{itemize}
% [Tutor Note: 原文 p.16 第二条把 Fukaya--Seidel 范畴印作“Fukaya-Siedel”（页图可辨，Siedel 系 Seidel 的拼写笔误）；人名按惯例还原为 Seidel，译文不另作注。]

由紧复曲面的分类，del Pezzo 曲面要么是 $\PP^2$ 在 $9-d$ 个点处的爆破（$0\le d\le8$），要么是 $\PP^1\times\PP^1$。因此，带光滑反典范除子的 del Pezzo 曲面配对总共有十个形变族。另一方面，带 $I_d$ 纤维的有理椭圆曲面也有十个形变族：对每个 $d$（$1\le d\le9$）各一个，另外对 $d=8$ 还有一个额外的族。于是当 $d\neq8$ 时，把次数为 $d$ 的 del Pezzo 曲面的每个形变族与带 $I_d$ 纤维的有理椭圆曲面的形变族对应起来并不难。对 $d=8$ 的两个族，还可以在拓扑上把它们等同起来。事实上，$\PP^1\times\PP^1$（或 Hirzebruch 曲面 $\mathbb{F}_1$）中一个光滑反典范除子的补集，与有理椭圆曲面中 $I_8$ 纤维的补集是微分同胚的，而该有理椭圆曲面与 $\PP^1\times\PP^1$（或相应地 $\mathbb{F}_1$）的超势紧化形变等价 [11, p.54]。

给定特殊拉格朗日纤维化 $X\to B$，在判别轨迹的补集上存在两个整仿射结构 [33]，分别由辛结构与复结构诱导。在 Calabi--Yau 曲面的情形，这两个整仿射结构可以如下定义：固定参考点 $u_0\in B_0$，其中 $B_0$ 是底 $B$ 中判别轨迹的补集。对每个 $\gamma\in H_1(L_{u_0},\Z)$ 与 $u_0$ 附近的 $u\in B_0$，选取连接 $u_0$ 与 $u$ 的道路 $\varphi$，并定义 $C_{\gamma,\varphi}$ 为 $\varphi$ 的像上的一个 $S^1$-丛，其纤维通过沿 $\varphi$ 的平行移动与 $\gamma$ 同伦。于是辛仿射坐标定义为
\begin{equation}
a_\gamma(u):=\int_{C_{\gamma,\varphi}}\omega, \tag{6.6.1}
\end{equation}
它因特殊拉格朗日条件而是良定义的。直接计算表明，不同参考点 $u_0$ 与道路 $\varphi$ 的选取之间的过渡函数属于 $\operatorname{GL}(2,\Z)\rtimes\R^2$。若把式 (6.6.1) 中的 $\omega$ 换成 $\operatorname{Im}\Omega$，便得到复仿射坐标。
% [Tutor Note: (6.1) 里环道记号是 $C_{\gamma,\varphi}$，其中第二个下标是道路 $\varphi$（varphi），与超势 $W$ 无关；全节的道路一律用 $\varphi$，不要与后文 $\Phi$（自同构）、$\phi$（单位周期映射）混同。]
